\section{Discussion}
\label{sec:discussion}

\subsection{Verdict against the criteria fixed in advance}

\S\ref{sec:setup} committed to three tests before any run, and E1 answers all three against
the proposal.

\begin{ttkey}
\lead{1. Does depth beat width at a matched clause budget?} Not on CIFAR-10. The \mctm
reaches \accMctm{}\% against \accSingleMatched{}\% for a single layer of the same total size
and \accSingleRf{}\% for a single layer with the stack's $9\times9$ receptive field. It does
not even reach the \accSingleLOne{}\% its own frozen layer~1 achieves alone: the second layer
\emph{removes} accuracy that was already there.

\lead{2. Does the spatial map earn its cost?} On CIFAR-10 the question barely arises --- both
stacked variants sit far below the single-layer baselines. On \code{near}
(\S\ref{sec:synthetic}) the answer is an emphatic yes, and it is the clearest positive result
in this report.

\lead{3. Does the greedy supervised objective buy anything?} No --- it is actively harmful.
Replacing the trained layer~1 with \emph{random} clauses raises the stack from \accMctm{}\%
to \accMctmRandom{}\%. Random three-literal clauses fire on \fireMctmRandom{}\% of positions
against \fireMctm{}\% for trained ones, and that density difference alone outweighs
everything the supervised objective learned.
\end{ttkey}

\subsection{What is actually broken}

The two risks identified before the experiments both materialised, and they compound.

\lead{Density collapse is the proximate cause.} Training layer~1 to classify rewards specific
clauses, which fire rarely; E0 shows every gain in density costs layer-1 accuracy, and the
most accurate setting fires on \densitySparsestFire{}\% of positions. Layer~2 then faces an
input where positive literals are almost never satisfied and negations almost always are,
while Type~I feedback rewards including whatever is $1$ in the drawn patch.
Table~\ref{tab:comp} shows where that ends: \negMctm{}\% of the literals in an \mctm layer-2
clause are negations. The second layer does not learn compositions of features. It learns
that things are absent.

\lead{Greedy saturation is the deeper cause.} E2 separates architecture from training. Given
features that are actually good, the stack clears a ceiling a single layer provably cannot
($\synXorFlatStackOracleSixFour{}\%$ against $75\%$, at a matched channel count). Given
features produced by the greedy scheme, it does not ($\synXorMctmTask{}\%$). The architecture
is not what fails.

\begin{ttcallout}[What survives]
Two things are worth keeping from the proposal. First, stacking convolutional Tsetlin layers
is genuinely more expressive than one layer --- demonstrated, not argued, on a task whose
one-layer ceiling is known. Second, \emph{keeping the spatial map} rather than globally
OR-pooling it is what makes relational concepts reachable at all: on \code{near} the spatial
stack reaches $\synNearMctmTask{}\%$ while the global-OR stack is pinned at chance whatever
features it is given. The proposal's distinctive architectural choice is the right one. Its
training scheme is not.
\end{ttcallout}


\subsection{What would have to change}

The greedy supervised scheme in the proposal couples two things that want opposite settings.
Layer~1 is trained to classify, which rewards \emph{specific} clauses --- many included
literals, rare firing --- and layer~2 needs \emph{general} clauses whose activations are
common enough to be conjoined. E0 (\S\ref{sec:density}) measures that coupling directly:
across the sweep, every increase in layer-1 density costs layer-1 accuracy. A design that
wants both cannot get them from one objective.

Three directions follow, in increasing order of how much they change the proposal.

\begin{itemize}[leftmargin=1.2em]
\item \lead{Give layer 1 a different objective.} Nothing forces a feature layer to be trained
  by classification. A reconstruction objective --- clauses that predict held-out features of
  a patch from the rest, as in the Tsetlin machine autoencoder --- preserves information
  rather than discarding whatever did not help the local classifier, and it has no reason to
  drive clauses towards extreme specificity. This is the direct analogue of the unsupervised
  greedy pre-training that made layer-wise deep learning work in the first place.
\item \lead{Calibrate density explicitly.} Rather than hoping $s$ lands somewhere usable, fix
  a target per-channel firing rate and adjust each clause towards it --- a Boolean analogue of
  normalisation. The measurement in \S\ref{sec:density} is the diagnostic such a scheme would
  optimise.
\item \lead{Propagate credit instead of freezing.} The interesting question the proposal
  defers (its \S23, Q9) is whether feedback can reach a lower layer at all. Clause inclusion
  is a usable credit path: if a layer-2 clause receives Type~I feedback and includes the
  literal ``channel $k$ active at offset $(dy,dx)$'', then layer-1 clause $k$ can receive
  Type~I feedback at that position. That is a discrete, rule-based analogue of
  backpropagation, it requires no gradients, and it removes the greedy-saturation problem by
  construction. It is a bigger change than anything tested here, and it is the part of the
  idea most worth pursuing.
\end{itemize}

\subsection{Limitations of this evaluation}

\lead{Scale.} These are small Tsetlin machines by CIFAR-10 standards --- hundreds of clauses
where the literature uses thousands --- chosen so that ten arms across three seeds fit in a
few GPU-hours. Absolute accuracies are correspondingly modest and should not be read as
CIFAR-10 results for Tsetlin machines. The comparison is between arms at a matched clause
budget, and that is what the conclusions rest on.

\lead{One stack depth, and a gap between the two datasets.} Only two-layer stacks were run.
E2 shows CIFAR-10 was, in part, the wrong place to look: its classes are not obviously
concepts that need a conjunction or a relation between separated parts, which is the only
thing the second layer supplies. That makes the CIFAR-10 result a statement about the greedy
training scheme on a natural-image task, not about the architecture --- E2 is what speaks to
the architecture. Whether a relational advantage of the kind \code{near} isolates ever shows
up on natural images is untested here, and is the obvious next question.

\lead{The oracle is an upper bound, not a method.} \code{build\_oracle\_l1} writes the right
features in by hand. It establishes that the architecture can use good features, and
therefore that training is the binding constraint; it does not suggest how to obtain such
features without knowing the answer. The three directions above are proposals, not results.

\lead{Hyper-parameters.} $T$ was fixed at the clause count and not tuned per arm, and layer-2
specificity was tied to layer-1's. A wider search could move the numbers; it is unlikely to
close the gaps reported in \S\ref{sec:results}, which are large relative to seed spread.
